% !TeX root = linear-codes.tex
\section{Circuits and walks}
\label{sec:approach}

Put $w=d-1$.
Fix $\eps>0$.  We work in the range
$$
\ln n\ll w\le(1/2-\eps)n.
$$
Define the redundancy parameter
$$
R\defeq\ln\frac{2^n}{V}.
$$
Our assumptions imply $R=\Theta(n)$.  Fix
$0<\zeta<1$, and choose the integer $m$ so that, with
$$
\tau\defeq V2^{-m},
$$
we have
$$
\frac{(1-\zeta)R}{2}
\le\tau<(1-\zeta)R.
$$
Such an $m$ exists because this interval has multiplicative width
two.

Choose $M$ uniformly from $\field_2^{m\times n}$ and put
$$
\kero\defeq\setbuilder{x\in\ball}{Mx=0}.
$$
A \emph{circuit} in $\ball$ is an inclusion-minimal linearly
dependent subset.  If $\circuit_\ell$ is the collection of
$\ell$-element circuits in $\ball$, then every
$S\in\circuit_\ell$ satisfies
$$
\sum_{x\in S}x=0,
\qquad
\rank(S)=\ell-1,
$$
where $3\le\ell\le n+1$.

Define the probability kernel
$$
\mu(x)\defeq\frac{\indicate{x\in\ball}}{V}.
$$
Let $X_1,X_2,\ldots$ be independent with law $\mu$, and set
$$
p_\ell\defeq\Pr[X_1+\cdots+X_\ell=0].
$$
Every circuit in $\circuit_\ell$ gives $\ell!$ distinct ordered
tuples counted by this event.  Hence
$$
\size{\circuit_\ell}
\le\frac{V^\ell}{\ell!}p_\ell.
$$

Let $Z$ be the number of odd circuits contained in $\kero$.
Since the span of an $\ell$-element circuit has dimension
$\ell-1$,
$$
\Pr[S\subseteq\ker M]
=2^{-m(\ell-1)}
=\nparen{\frac{\tau}{V}}^{\ell-1}.
$$
Consequently,
$$
\begin{aligned}
\E Z
&=\sum_{\substack{3\le\ell\le n+1\\\ell\text{ odd}}}
  \size{\circuit_\ell}
  \left(\frac{\tau}{V}\right)^{\ell-1}\\
&\le\frac V\tau
  \sum_{\substack{\ell\ge3\\\ell\text{ odd}}}
  \frac{\tau^\ell}{\ell!}p_\ell.
\end{aligned}
$$

It suffices to find a choice of $M$ for which $Z=0$.  Indeed,
every zero-sum subset can be partitioned into circuits.  Hence, if
$\kero$ contains no odd circuit, then every dependency among its
elements has even cardinality.  Let $\mathbf K$ be the matrix whose
columns are the elements of $\kero$.  The preceding statement means
that
$$
\mathbf K a=0
\quad\Longrightarrow\quad
\mathbf1^\transpose a=0.
$$
Thus $\mathbf1^\transpose$ lies in the row space of $\mathbf K$,
and there is a row vector $v$ satisfying
$$
v\mathbf K=\mathbf1^\transpose.
$$
Appending $v$ to $M$ then produces a parity-check matrix for a
distance-$d$ code.

We now express the sum above as a transition probability.  Let
$J_t$ be Poisson with mean $t$, independently of the $X_i$, and
define the continuous-time random walk on $\field_2^{n+1}$ by
$$
W_t\defeq\sum_{i=1}^{J_t}(X_i,1).
$$
The walk starts at zero.  Write
$$
e_{n+1}\defeq(0,\ldots,0,1).
$$
Conditioning on the number of steps gives
$$
e^\tau\Pr[W_\tau=e_{n+1}]
=
\sum_{\substack{\ell\ge3\\\ell\text{ odd}}}
\frac{\tau^\ell}{\ell!}p_\ell,
$$
where the one-step term vanishes because $0\notin\ball$.
Consequently,
$$
\E Z
\le
\frac{Ve^\tau}{\tau}\Pr[W_\tau=e_{n+1}].
$$
It remains to prove
$$
\Pr[W_\tau=e_{n+1}]
<\frac{\tau e^{-\tau}}{V}.
$$
